\clearpage
\markboth{参考文献}{参考文献}
\phantomsection
\addcontentsline{toc}{section}{参考文献}
\begingroup
    \centering
    \zihao{-3}\heiti\bfseries 参考文献
    \par
    \vspace{\baselineskip}
\endgroup

\begin{enumerate}[label={[\arabic*]}, leftmargin=2em, itemsep=0pt, parsep=0pt]
    \zihao{-4}\songti
    \item 代世峰，任德贻，唐跃刚，等．高硫煤中硫的地质演化模式——以内蒙古乌达矿区为例[J]．地质论评，2001，47(4)：383-387.
    \item Kostova I, Petrov O, Kortenski J. Mineralogy, geochemistry and pyrite content of Bulgarian subbituminous coals, Pernik Basin[J]. Geological Society, London, Special Publications, 1996, 109(1): 301-314.
    \item 叶连俊．生物有机机质成矿作用和成矿背景[M]．北京：海洋出版社，1998: 75-84.
    \item Stach E, Mackowsky M, Teichmuller M, et al. Stach's Textbook of Coal Petrology[J]. Borntraeger, Stuttgart, 1982: 12-89.
    \item 代世峰，马风学．乌达矿区高硫煤中菌藻类体的发现及意义[J]．中国矿业大学学报，1999，28(1)：57-60.
    \item Chou C L. Geologic factors affecting the abundance, distribution, and speciation of sulfur in coals[C]. //Geology of Fossil Fuels, Proc 30th Int Geol Congress. 1997: 47-57.
    \item 郑宝山．地方性氟中毒及工业氟污染研究[M]．北京：中国环境科学出版社，1992：169-172.
    \item 代世峰．煤中伴生元素的地质地球化学习性与富集模式[D]．北京：中国矿业大学(北京)，2002.
    
    \item[\dots\dots]
    
    \item[ [23] ] 于潇，刘义，柴跃廷，等．互联网药品可信交易环境中主体资质审核备案模式[J]．清华大学学报(自然科学版)，2012，52(11)：1518-1523.
    
    \item[\dots\dots]
    
    \item[ [220] ] 哈里森，沃尔德伦．经济数学与金融数学[M]．谢远涛，译．北京：中国人民大学出版社，2012：235-236.
\end{enumerate}

\clearpage
\mbox{}\clearpage 
